\documentclass[UTF8,zihao=-4]{ctexart}
\usepackage[a4paper,margin=2.2cm]{geometry}
\usepackage{array}
\usepackage{longtable}
\usepackage{hyperref}
\usepackage{verbatim}
\hypersetup{hidelinks}
\setlength{\parindent}{0pt}
\setlength{\parskip}{0.55em}
\setcounter{secnumdepth}{0}
\renewcommand{\arraystretch}{1.25}
\emergencystretch=4em
\sloppy
\title{05 设计模式与软件构造复习讲义}
\author{}
\date{}
\begin{document}
\maketitle
\tableofcontents
\newpage


\begin{quote}
范围：\texttt{设计模式.pdf}、\texttt{代码设计.pdf}、\texttt{软件构造.pdf}；结合 Exam/Review 中策略模式、抽象工厂、代码质量、重构、审查、测试和构造管理内容。
结构：先整理 PPT 主线，再放对应考试模板和快速背诵清单。
\end{quote}

\bigskip\hrule\bigskip

\section{5. 设计模式}

\subsection{5.1 设计模式的构成}

一个设计模式通常包含：

\begin{itemize}
\item 名称
\item 问题
\item 解决方案
\item 应用场景
\item 注意事项
\end{itemize}

模式不是为了“套模板”，而是把稳定抽象和变化点分开。

\subsection{5.2 Strategy 策略模式}

定义：

\begin{quote}
定义一组算法，把每个算法封装起来，并让它们可以互换，使算法变化独立于使用算法的客户。
\end{quote}

结构：

\begin{itemize}
\item Context：持有 Strategy。
\item Strategy：算法接口。
\item ConcreteStrategy：具体算法。
\end{itemize}

适用场景：

\begin{itemize}
\item 同一业务存在多种算法。
\item 需要运行时切换算法。
\item 代码中存在大量按类型分支选择行为。
\end{itemize}

例子：充值手续费策略

\small
\begin{verbatim}
public interface FeeStrategy {
    BigDecimal calculateFee(BigDecimal amount);
}

public class BankCardFeeStrategy implements FeeStrategy {
    public BigDecimal calculateFee(BigDecimal amount) {
        return amount.multiply(new BigDecimal("0.005"));
    }
}

public class AlipayFeeStrategy implements FeeStrategy {
    public BigDecimal calculateFee(BigDecimal amount) {
        return BigDecimal.ZERO;
    }
}
\end{verbatim}
\normalsize

Context：

\small
\begin{verbatim}
public class RechargeServiceImpl implements RechargeService {
    private final FeeStrategy feeStrategy;

    public RechargeServiceImpl(FeeStrategy feeStrategy) {
        this.feeStrategy = feeStrategy;
    }

    public void recharge(BigDecimal amount) {
        BigDecimal fee = feeStrategy.calculateFee(amount);
        // 执行业务流程
    }
}
\end{verbatim}
\normalsize

优点：

\begin{itemize}
\item 消除复杂条件分支。
\item 新增算法时新增策略类。
\item 符合 OCP。
\end{itemize}

代价：

\begin{itemize}
\item 类数量增加。
\item 客户端需要选择策略。
\item Strategy 与 Context 之间需要约定输入输出。
\end{itemize}

\subsection{5.3 Abstract Factory 抽象工厂}

定义：

\begin{quote}
提供一个创建一组相关或相互依赖对象的接口，而不指定它们的具体类。
\end{quote}

结构：

\begin{itemize}
\item AbstractFactory
\item ConcreteFactory
\item AbstractProduct
\item ConcreteProduct
\item Client
\end{itemize}

适用场景：

\begin{itemize}
\item 系统需要创建一组相关产品。
\item 产品族需要整体替换。
\item 客户端不应依赖具体产品类。
\end{itemize}

例子：不同平台 UI 组件族

\small
\begin{verbatim}
public interface UIFactory {
    Button createButton();
    Dialog createDialog();
}

public class WebUIFactory implements UIFactory {
    public Button createButton() {
        return new WebButton();
    }

    public Dialog createDialog() {
        return new WebDialog();
    }
}
\end{verbatim}
\normalsize

优点：

\begin{itemize}
\item 保证同一产品族一起使用。
\item 客户端依赖抽象。
\item 切换产品族集中在工厂。
\end{itemize}

缺点：

\begin{itemize}
\item 新增产品种类会修改工厂接口。
\item 产品族稳定时效果最好。
\end{itemize}

\subsection{5.4 Factory Method 工厂方法}

定义：

\begin{quote}
父类定义创建对象的接口，由子类决定实例化哪一个具体类。
\end{quote}

适用场景：

\begin{itemize}
\item 父类无法提前确定要创建的具体对象。
\item 创建逻辑要延迟到子类。
\end{itemize}

与抽象工厂区别：

\begin{itemize}
\item 工厂方法关注创建一个产品。
\item 抽象工厂关注创建一组相关产品。
\end{itemize}

\subsection{5.5 Singleton 单例}

定义：

\begin{quote}
保证一个类只有一个实例，并提供全局访问点。
\end{quote}

使用时要谨慎：

\begin{itemize}
\item 单例容易形成全局状态。
\item 会增加隐式耦合。
\item 测试替换不方便。
\end{itemize}

适合：

\begin{itemize}
\item 无状态工具入口。
\item 全局唯一资源管理器。
\item 配置对象。
\end{itemize}

不适合：

\begin{itemize}
\item 承载业务状态。
\item 替代依赖注入。
\end{itemize}

\subsection{5.6 Iterator 迭代器}

定义：

\begin{quote}
在不暴露聚合对象内部表示的前提下，顺序访问其中元素。
\end{quote}

结构：

\begin{itemize}
\item Iterator
\item ConcreteIterator
\item Aggregate
\item ConcreteAggregate
\end{itemize}

优点：

\begin{itemize}
\item 隐藏集合内部结构。
\item 简化聚合对象接口。
\item 支持多种遍历方式。
\item 支持多个遍历同时进行。
\end{itemize}

例子：

客户端使用迭代器遍历充值记录，不关心底层是数组、链表、数据库游标还是分页结果。

\bigskip\hrule\bigskip

\bigskip\hrule\bigskip

\section{6. 代码设计与软件构造}

\subsection{6.1 代码设计关注点}

代码设计关注可读性、可维护性和可验证性。

主要内容：

\begin{itemize}
\item 布局
\item 命名
\item 注释
\item Javadoc
\item 小任务拆分
\item 复杂决策表达
\item 数据使用
\item 显式依赖
\item 断言
\end{itemize}

\subsection{6.2 布局}

原则：

\begin{itemize}
\item 缩进表达逻辑层级。
\item 空行分隔不同语义块。
\item 长表达式拆行。
\item 相似语句对齐。
\item 复杂控制结构保持清楚。
\end{itemize}

坏布局会让正确代码难以审查。

\subsection{6.3 命名}

命名要求：

\begin{itemize}
\item 表达含义。
\item 与领域术语一致。
\item 避免误导。
\item 避免无意义缩写。
\item 避免过长。
\end{itemize}

Java 约定：

\begin{itemize}
\item 包名：小写单词。
\item 类名和接口名：PascalCase。
\item 变量和方法：camelCase。
\item 常量：大写字母加下划线。
\end{itemize}

语义规则：

\begin{itemize}
\item 类、属性、数据对象用名词。
\item 接口可用名词或形容词。
\item 方法用动词或动词加名词。
\end{itemize}

示例：

\begin{itemize}
\item \texttt{RechargeService}
\item \texttt{CampusCardRepository}
\item \texttt{findRechargeRecordsByCardNo}
\item \texttt{MAX\_RECHARGE\_AMOUNT}
\end{itemize}

\subsection{6.4 注释与 Javadoc}

Java 注释：

\begin{itemize}
\item \texttt{//}：单行注释。
\item \texttt{/* */}：普通块注释。
\item \texttt{/** */}：文档注释。
\end{itemize}

Javadoc 常用标签：

\begin{itemize}
\item \texttt{@author}
\item \texttt{@version}
\item \texttt{@see}
\item \texttt{@since}
\item \texttt{@deprecated}
\item \texttt{@param}
\item \texttt{@return}
\item \texttt{@throws}
\end{itemize}

注释原则：

\begin{itemize}
\item 解释意图、约束和不直观决策。
\item 不重复代码表面含义。
\item 公共 API 要说明参数、返回值和异常。
\end{itemize}

\subsection{6.5 AI 代码生成 Prompt}

四个元素：

\begin{itemize}
\item Context：上下文。
\item Goal：目标。
\item Format：输出格式。
\item Steps：步骤。
\end{itemize}

示例：

\small
\begin{verbatim}
Context:
项目使用 Spring Boot，分层为 controller/service/repository。

Goal:
实现校园卡充值服务接口和实现类。

Format:
输出 Java 接口、实现类和单元测试。

Steps:
先定义 DTO，再定义 Service 接口，再写实现逻辑，最后写测试。
\end{verbatim}
\normalsize

\subsection{6.6 软件构造}

SWEBOK 中软件构造覆盖：

\begin{itemize}
\item 编码
\item 验证
\item 单元测试
\item 集成测试
\item 调试
\end{itemize}

构造不是“写代码”一个动作，而是把设计转化为可运行、可验证、可维护的软件制品。

\subsection{6.7 构造技术}

主要技术：

\begin{itemize}
\item 可理解的源代码
\item 类、枚举、变量、常量
\item 控制结构
\item 错误处理
\item 代码级安全
\item 单元测试
\item 集成测试
\item 调试
\end{itemize}

\subsection{6.8 代码审查}

课程给出的审查经验：

\begin{itemize}
\item 每次审查少于 200 到 400 行代码。
\item 审查速度控制在每小时 300 到 500 行。
\item 单次审查时间控制在 60 到 90 分钟。
\item 开发者先对代码做注解。
\item 使用检查表。
\item 确认缺陷已修复。
\item 保持积极的审查文化。
\item 借助工具支持轻量化审查。
\end{itemize}

AI 辅助代码审查适合：

\begin{itemize}
\item 识别坏味道。
\item 检查安全问题。
\item 检查风格问题。
\item 发现遗漏测试。
\end{itemize}

最终结论要由人结合需求、设计和上下文确认。

\subsection{6.9 集成与构建}

集成策略：

\begin{itemize}
\item 大爆炸集成：一次性集成所有模块。
\item 增量集成：逐步集成并验证模块。
\end{itemize}

增量集成更利于定位问题和控制风险。

构建：

\begin{quote}
将源代码、资源、配置和依赖转化为可执行或可部署制品。
\end{quote}

构建依赖配置管理：

\begin{itemize}
\item 版本一致。
\item 依赖可追踪。
\item 构建脚本受控。
\item 发布制品可复现。
\end{itemize}

\subsection{6.10 构造管理}

构造管理内容：

\begin{itemize}
\item 分配构造任务。
\item 跟踪代码进度。
\item 度量复杂度。
\item 度量代码行数。
\item 度量注释比例。
\item 管理分支、合并和发布。
\item 执行代码审查和测试门禁。
\end{itemize}

\subsection{6.11 坏味道与重构}

重构定义：

\begin{quote}
在不改变软件外部行为的前提下，改善内部结构。
\end{quote}

常见坏味道：

\begin{itemize}
\item 重复代码。
\item 长方法。
\item 大类。
\item 长参数列表。
\item Magic Number 或 Magic String。
\item 空 catch。
\item 过多注释。
\item 某个类过度访问其他类属性。
\item Switch 或 if-else 按类型分发行为。
\end{itemize}

对应改进：

\begin{longtable}{|>{\raggedright\arraybackslash}p{0.24\textwidth}|>{\raggedright\arraybackslash}p{0.31\textwidth}|>{\raggedright\arraybackslash}p{0.31\textwidth}|}
\hline
\textbf{坏味道} & \textbf{问题} & \textbf{重构} \\
\hline
\endfirsthead
\hline
\textbf{坏味道} & \textbf{问题} & \textbf{重构} \\
\hline
\endhead
重复代码 & 隐式耦合 & 提取方法或提取类 \\
\hline
长方法 & 难理解、难测试 & 提取方法 \\
\hline
大类 & 职责过多 & 提取类、按职责拆分 \\
\hline
长参数列表 & 调用复杂 & 引入参数对象 \\
\hline
Magic Number & 含义不清 & 提取常量 \\
\hline
空 catch & 吞掉错误 & 记录日志或抛出业务异常 \\
\hline
过多注释 & 结构表达力不足 & 改名、拆分、重组逻辑 \\
\hline
过度访问其他类属性 & 职责位置错误 & 移动方法或委托 \\
\hline
类型分支 & 违反 OCP & 策略模式或多态 \\
\hline
\end{longtable}

\subsection{6.12 TDD}

TDD 循环：

\begin{enumerate}
\item Red：先写失败测试。
\item Green：写最少代码让测试通过。
\item Refactor：在测试保护下重构。
\end{enumerate}

适合：

\begin{itemize}
\item 规则明确的业务逻辑。
\item 算法。
\item 服务层。
\item 状态迁移。
\end{itemize}

\subsection{6.13 AI 编程原则}

使用 AI 写代码时：

\begin{itemize}
\item 明确意图。
\item 逐步细化。
\item 保持对代码的理解。
\item 先定义测试或验收标准。
\item 经过验证后再提交。
\end{itemize}

不能把 AI 输出直接当作正确实现。需要用需求、设计、测试和代码审查约束它。

\bigskip\hrule\bigskip

\bigskip\hrule\bigskip

\section{对应考试模板}

\subsection{7.9 设计模式题模板}

题目出现大量按类型选择算法：

\small
\begin{verbatim}
if (payType.equals("BANK")) { ... }
else if (payType.equals("ALIPAY")) { ... }
else if (payType.equals("WECHAT")) { ... }
\end{verbatim}
\normalsize

答：

\begin{itemize}
\item 使用 Strategy。
\item 抽象 \texttt{PaymentStrategy}。
\item 每种支付方式一个具体策略。
\item \texttt{RechargeService} 依赖 \texttt{PaymentStrategy} 或策略工厂。
\item 新增支付方式时新增策略类，减少修改原业务流程。
\end{itemize}

题目出现一组相关产品整体替换：

\begin{itemize}
\item 使用 Abstract Factory。
\item 抽象产品和抽象工厂。
\item 具体工厂创建同一产品族。
\end{itemize}

题目要求隐藏集合结构：

\begin{itemize}
\item 使用 Iterator。
\item 客户端通过迭代器访问元素。
\end{itemize}

\subsection{7.10 代码质量题模板}

答题结构：

\small
\begin{verbatim}
问题：
违反原则：
影响：
重构：
测试：
\end{verbatim}
\normalsize

示例：

\small
\begin{verbatim}
问题：充值方法过长，同时处理参数校验、支付调用、余额更新、记录保存和异常处理。
违反原则：SRP、SOFA。
影响：代码难理解、难测试、修改支付逻辑会影响余额更新逻辑。
重构：提取校验方法、支付网关、余额更新方法和记录保存方法；支付方式变化使用 Strategy。
测试：为正常充值、支付失败、校园卡挂失、金额非法分别编写单元测试。
\end{verbatim}
\normalsize

\bigskip\hrule\bigskip

\section{快速背诵清单}

\subsection{设计模式}

\begin{itemize}
\item Strategy：封装可互换算法。
\item Abstract Factory：创建一组相关产品。
\item Factory Method：子类决定创建哪个产品。
\item Singleton：唯一实例和全局访问点。
\item Iterator：隐藏集合内部结构并顺序访问元素。
\end{itemize}

\subsection{构造}

\begin{itemize}
\item 构造 = 编码 + 验证 + 单元测试 + 集成测试 + 调试。
\item 代码审查要小批量、限时、用检查表。
\item 重构不改变外部行为，只改善内部结构。
\item TDD：Red、Green、Refactor。
\end{itemize}


\end{document}
